The following discussion is presented as general background knowledge: no external evidence excerpts were provided to support specific empirical claims, and the statements below should be interpreted as design recommendations and hypotheses rather than verified results. (General background knowledge.)

We recommend enforcing a strict JSON schema per section (fields for rendered text, cited evidence identifiers, and provenance metadata) so each authored unit yields a compact, verifiable record that supports automated schema/citation checks and simplifies human inspection.

Concretely, require claim-level evidence_ids (automated citation–assertion checking); capture fine-grained provenance for each snippet (source id, retrieval query/score/rank, offsets, transformations); package versioned pipelines and configuration (indexes, prompt templates, schema versions) for reproducibility; and broaden evaluation (more domains and structured human assessments, plus metrics over the JSON outputs).

Limitations include dependence on retrieval quality, possible rigidity of schemas, and the need to balance automated blocking with human-in-the-loop review; adopting JSON-constrained outputs and explicit provenance can improve auditability in principle but requires implementing these components and empirical validation. (General background knowledge.)